\section{Methods}
\label{sec:methods}

% All parameter values in this section were read from the cited source files during the
% 2026-08-07 audit. See manuscript/audit/SOURCE_OF_TRUTH.md for the full inventory and
% manuscript/audit/DRAFT_STALENESS.md for the reconciliation against earlier drafts.

\subsection{Overview}
\label{sec:methods-overview}

The workflow separates a heterogeneous simulation frame into physically meaningful spatial
domains, represents each domain with a fixed structural descriptor, decides whether the
domain resembles anything in a known atlas, and---independently---asks which physical
reference structure, if any, is compatible with it. Detection and identification are
deliberately decoupled: the first is a statistical statement about a training distribution,
the second a comparison against externally derived reference structures. A domain may be
flagged as novel and then remain unidentified, which is a reportable outcome rather than a
failure.

\subsection{Gas removal and physical connectivity}
\label{sec:gas}

% SOURCE:
% results/motif_pipeline/frozen_gas_m3_pipeline_v1/gas_m3_pipeline.py
% function: radius_normalized_coord_r4_gas_screen
Particle separations are made dimensionless by the sum of radii,
\begin{equation}
  s_{ij} \;=\; \frac{d_{ij}}{R_i + R_j},
  \label{eq:sij}
\end{equation}
so that all screens are independent of the absolute colloid size. Two particles are treated
as belonging to the same condensed region when $s_{ij} < 1.62$; connected components under
this relation define the \emph{physical components} of the frame, and particles in no
component are labelled gas (component index $-1$). A separate, tighter cutoff
$s_{ij} < M_{\mathrm{graph}}$ with $M_{\mathrm{graph}} = 1.35$ defines the
\emph{contact graph} used for the coordination and exposure quantities of
\cref{sec:centres}. Both thresholds are frozen constants of the screen.

Physical connectivity and structural identity are distinct. A physical component is a
region of space in which particles are in contact; it may contain more than one crystal
structure, and two structurally identical grains may form separate components. The
component is the unit on which the characteristic length scale and the descriptor context
are defined, not the unit that carries a structural label.

\subsection{Reliable-centre selection}
\label{sec:centres}

% SOURCE:
% results/motif_pipeline/frozen_gas_m3_pipeline_v1/gas_m3_pipeline.py
% function: _exposure_and_coord
% selector name: phase_blind_coord_exposure_v1
Particles at a free surface see an incomplete neighbour shell, so their local descriptors
report the surface rather than the bulk structure. Two per-particle quantities are computed
over same-component contact neighbours ($s_{ij} < 1.35$): the contact coordination
$n_i$, and the \emph{exposure}
\begin{equation}
  e_i \;=\; \frac{1}{n_i}\,\Bigl\lVert \sum_{j} \widehat{(\bm{r}_j - \bm{r}_i)} \Bigr\rVert,
  \label{eq:exposure}
\end{equation}
the norm of the mean unit vector to the contact neighbours, set to $1.0$ for an isolated
particle. For an interior particle the neighbour directions cancel and $e_i \to 0$; for a
particle on a surface they do not, and $e_i$ approaches unity. A particle is excluded as a
descriptor centre when
\begin{equation}
  e_i > \tau_e \;\;\text{and}\;\; n_i < \tau_c,
  \qquad \tau_e = 0.14,\;\; \tau_c = 13 .
  \label{eq:m3}
\end{equation}

The conjunction matters: one-sidedness alone would remove well-coordinated particles at a
flat facet, and low coordination alone would remove interior particles in an open structure.

% SOURCE:
% results/motif_pipeline/soap_atlas_core3_phase2_v1/scripts/core3_pipeline.py
% function: soap_global_d0
This is a restriction on which particles are used as \emph{centres}, not a deletion of
particles. Excluded particles remain part of the physical structure and remain in the
neighbour context of every retained centre: the descriptor routine passes all non-gas
particles as context while evaluating the descriptor only at retained centres. We therefore
refer to the operation as reliable-centre selection rather than surface removal.

\subsection{Species-aware SOAP representation}
\label{sec:soap}

% SOURCE:
% results/motif_pipeline/soap_atlas_core3_phase2_v1/scripts/core3_pipeline.py
% constants: Z_LARGE=11, Z_SMALL=17, N_MAX=8, L_MAX=6, R_CUT_NORM=3.0, SIGMA_NORM=0.10
% function: _soap (DScribe SOAP)
Each retained centre is described by the smooth overlap of atomic positions. The neighbour
density of species $\alpha$ around centre $i$ is a sum of Gaussians,
\begin{equation}
  \rho^{\alpha}_{i}(\bm{r}) \;=\;
  \sum_{j \in \alpha,\; r_{ij} < r_{\mathrm{cut}}}
  \exp\!\left(-\frac{\lVert \bm{r} - \bm{r}_{ij} \rVert^{2}}{2\sigma^{2}}\right),
  \label{eq:soap-density}
\end{equation}
expanded in radial basis functions $g_n$ and spherical harmonics $Y_{lm}$,
\begin{equation}
  \rho^{\alpha}_{i}(\bm{r}) \;=\; \sum_{nlm} c^{\alpha,i}_{nlm}\, g_n(r)\, Y_{lm}(\hat{\bm{r}}),
  \label{eq:soap-expansion}
\end{equation}
and reduced to the rotationally invariant power spectrum
\begin{equation}
  p^{\alpha\beta,i}_{nn'l} \;=\; \pi\sqrt{\tfrac{8}{2l+1}}\;
  \sum_{m=-l}^{l} \bigl(c^{\alpha,i}_{nlm}\bigr)^{*} c^{\beta,i}_{n'lm}.
  \label{eq:soap-power}
\end{equation}

The representation is invariant to translation by construction, to rotation through
\cref{eq:soap-power}, and to permutation of identical particles because
\cref{eq:soap-density} sums over neighbours without ordering them. The Gaussian width
$\sigma$ smooths thermal displacement, so a finite-temperature realisation of a structure
maps near its ideal counterpart.

Two species channels are used, labelled by placeholder atomic numbers $Z = 11$ (LARGE) and
$Z = 17$ (SMALL), giving the species pairs $\{LL, LS, SS\}$. With $n_{\max} = 8$ and
$l_{\max} = 6$ and DScribe's crossover convention the descriptor length is
\begin{equation}
  \frac{(N_s n_{\max})(N_s n_{\max} + 1)}{2}\,(l_{\max}+1)
  \;=\; \frac{16 \cdot 17}{2}\cdot 7 \;=\; 952 ,
  \label{eq:soap-dim}
\end{equation}
with $N_s = 2$. The production code asserts this length at run time and raises if the
descriptor configuration changes.

% SOURCE:
% results/motif_pipeline/soap_atlas_core3_phase2_v1/scripts/core3_pipeline.py
% function: component_d0
The cutoff and width are set as multiples of a characteristic length rather than as absolute
distances:
\begin{equation}
  r_{\mathrm{cut}} = 3.0\, d_0, \qquad \sigma = 0.10\, d_0 ,
\end{equation}
where $d_0$ is the \emph{median nearest-neighbour distance within the physical component}.
The cutoff therefore spans three neighbour shells, and the descriptor of a structure is
unchanged if all colloid diameters and the box are rescaled together. Descriptor vectors are
$L_2$-normalised row-wise before any further processing.

% SOURCE:
% results/motif_pipeline/core3_comparator_benchmark_v1/scripts/development_scoring_common.py
% function: load_raw

\subsection{Species-conditioned standardisation and projection}
\label{sec:pca}

% SOURCE:
% results/motif_pipeline/core3_atlas_v3/scripts/v3_common.py
% functions: stratum_sample, fit_transform
The two species occupy systematically different local environments, so descriptors are
standardised \emph{per species}: for a centre of species $s$,
\begin{equation}
  \tilde{x}_{ij} \;=\; \frac{\hat{x}_{ij} - \mu_{s,j}}{\sigma_{s,j}},
  \qquad \hat{\bm{x}}_i = \bm{x}_i / \lVert \bm{x}_i \rVert_2 ,
  \label{eq:standardise}
\end{equation}
with $\mu_s$ and $\sigma_s$ estimated on training centres of that species only. A single
principal-component projection, shared across species, then maps the standardised
descriptors to
\begin{equation}
  \bm{z}_i \;=\; W^{\top} \tilde{\bm{x}}_i \in \mathbb{R}^{16}.
  \label{eq:pca}
\end{equation}
The projection returns $16$ learned linear combinations of the $952$ descriptor components;
it does not select $16$ of them.

The sample on which $W$ is fitted is balanced by construction. Each
(family, condition-group, species) stratum contributes exactly $128$ centres, split evenly
across the regions within that stratum, so a condition group with many trajectories does not
dominate the projection. Both $\mu_s, \sigma_s$ and $W$ are fitted on training data only.

\subsection{Domain-level fingerprint}
\label{sec:fingerprint}

% SOURCE:
% results/motif_pipeline/core3_atlas_v3/scripts/v3_common.py
% function: summary
Latent vectors are pooled to the domain level separately by species and concatenated,
\begin{equation}
  \bm{h} \;=\; \bigl[\, \bar{\bm{z}}_{L} \;\Vert\; \bar{\bm{z}}_{S} \,\bigr],
  \qquad
  \bar{\bm{z}}_{c} = \frac{1}{|C_c|}\sum_{i \in C_c} \bm{z}_i ,
  \label{eq:fingerprint}
\end{equation}
giving a $32$-dimensional fingerprint. Pooling separately by species preserves information
that a single mean would destroy: two structures can share an overall environment
distribution while differing in which sublattice occupies which environment.

Mean pooling has a known cost, quantified in \cref{sec:results}: a minority of novel
environments inside an otherwise familiar domain is averaged towards the familiar mean.

A domain is scored only when it satisfies a support requirement of at least $20$ retained
centres with at least $8$ per species; otherwise it is reported as unscoreable rather than
assigned a score.
% SOURCE:
% results/motif_pipeline/core3_atlas_v3/scripts/v3_01_build_shards.py
% constants: MIN_CENTERS=20, MIN_PER_SPECIES=8

\subsection{Open-set novelty detection}
\label{sec:ood}

% SOURCE:
% results/motif_pipeline/core3_atlas_v3/scripts/v3_common.py
% functions: fit_covariance, mahal, threshold
Known structure is modelled as a single distribution over domain fingerprints, with mean
$\bm{\mu}$ and covariance $\Sigma$ estimated on training domains. The novelty score is the
Mahalanobis distance
\begin{equation}
  \Dmah(\bm{h}) \;=\; \sqrt{ (\bm{h}-\bm{\mu})^{\top} \Sigma^{-1} (\bm{h}-\bm{\mu}) } .
  \label{eq:mahalanobis}
\end{equation}
The covariance is estimated with Ledoit--Wolf shrinkage,
\begin{equation}
  \hat{\Sigma} \;=\; (1-\shrink) S \;+\; \shrink \frac{\operatorname{tr}(S)}{p} I ,
  \label{eq:ledoitwolf}
\end{equation}
with $\shrink$ chosen analytically. The estimator is deterministic given the training rows.

% RESULT SOURCE:
% results/motif_pipeline/core3_atlas_v3/data/v3_threshold_diagnostic.csv
% columns: shrinkage
We note that in the present configuration the shrinkage is small---$0.0062$ and $0.0059$ for
the two training atlases of \cref{sec:results}---because $272$--$280$ group-balanced rows in
$32$ dimensions is a well-posed estimation problem. Shrinkage here is therefore a safeguard
rather than a necessity, in contrast to a regime where the covariance would be singular.

Every training condition group contributes exactly four fingerprints to \cref{eq:ledoitwolf},
by deterministic cyclic tiling of the group's regions. This gives each condition group equal
weight regardless of how many trajectories it contains.

The threshold is set on held-out calibration data. Per-domain scores are averaged within each
calibration condition group, and
\begin{equation}
  \tau \;=\; Q_{0.95}\bigl(\{\,\overline{\Dmah}^{(g)}\,\}_{g \in \mathrm{calibration}}\bigr),
  \label{eq:threshold}
\end{equation}
using the ``higher'' quantile convention. Scores are reported normalised,
$\rho = \Dmah/\tau$, and a domain is called novel when $\rho > 1$. Averaging within groups
before taking the quantile prevents a heavily sampled condition from dominating the
threshold.

\subsection{Data roles and firewalls}
\label{sec:splits}

% SOURCE:
% results/motif_pipeline/core3_atlas_v3/scripts/v3_00_predeclare.py
% artifact: results/motif_pipeline/core3_atlas_v3/data/V3_PREDECLARATION.json
Condition groups---distinct combinations of screening length and volume fraction---are the
unit of splitting, so trajectories from the same physical condition never appear on both
sides of a split. The split, the fitting configuration and the evaluation metrics were fixed
and hashed before any model was fitted; the predeclaration digest is recorded with the
artifacts.

Six condition groups are reserved as an untouched test set and were not opened during any
of the work reported here. Where provisional machine-generated labels are used, they are
confined to training: calibration and validation regions are entirely expert-labelled, and
this is asserted at run time rather than left as a convention.

\subsection{Strict unseen-family evaluation}
\label{sec:lofo}

% SOURCE:
% results/motif_pipeline/core3_atlas_v3/scripts/v3_03_compare.py
% function: m1_lofo
To test whether the representation supports families absent from training, each family in
turn is removed entirely: the standardisation, projection, covariance and threshold are all
re-fitted on the remaining families, and the removed family is scored as unseen. Positives
are the held-out family's calibration and validation domains; negatives are validation
domains of the retained families. We report AUROC, true-positive rate at $\tau$, and the
false-positive rate on known domains at the same $\tau$.

We distinguish this protocol, in which a \emph{fixed} descriptor is retained and only
statistical quantities are re-fitted, from strict retraining of a \emph{learned}
representation without the held-out family. The two answer different questions, and the
distinction is essential to any comparison against a supervised model.
\TODO{VERIFY: whether the GAT comparison in reports/gat\_pipeline retrains or only re-fits
statistics. This determines what can be claimed in \cref{sec:results}.}

\subsection{Physical reference matching}
\label{sec:matcher}

% SOURCE:
% results/statistical_ood_benchmark_v2/reference_conditioned_matcher_v1/build_and_run_matcher.py
Novelty detection returns no structural identity. A separate matcher compares a domain
against a library of references built from crystallographic information files, and is
permitted to abstain.

\paragraph{References.} Seven families are used. Each maps to a simulation run whose first
frame is the built realisation of the corresponding CIF; the matcher consumes these built
configurations rather than the CIF files directly. The ideal first frame serves as the
identity \emph{anchor} used for ranking. Eleven thermal frames per family, taken at fixed
intervals from an independent finite-temperature trajectory of the same family, define the
\emph{tolerance} ensemble used for acceptance. Anchors and tolerances therefore come from
different sources, and neither uses the candidate under test.

\paragraph{Domain extraction.} A candidate domain is extracted by coordination-based gas
rejection followed by selection of the largest connected condensed cluster, retaining every
particle in that cluster. A per-particle bond-order threshold is deliberately not used: it
is species-biased for asymmetric structures and distorts the composition on which the
species-resolved correlation functions depend.

\paragraph{Channels.} Three channels are compared. The species-resolved partial radial
distribution functions $g_{LL}, g_{LS}, g_{SS}$ are computed with a finite-domain
normalisation,
\begin{equation}
  g_{ab}(r) \;=\;
  \frac{N_{ab}(r)}{N^{\mathrm{ctr}}_{a}\, \rho_b\, 4\pi r^2 \Delta r\, S_a(r)} ,
  \label{eq:partial-rdf}
\end{equation}
% SOURCE:
% results/statistical_ood_benchmark_v2/matcher_v2_corrected_rdf/partial_rdf.py
where $N^{\mathrm{ctr}}_a$ is the number of species-$a$ centres, $\rho_b$ is the number
density of species $b$ estimated locally within the domain, and $S_a(r)$ is a
shell-containment correction: the fraction of the sphere of radius $r$ about a centre that
lies inside the domain, sampled over $26$ directions. Without $S_a(r)$ a finite crystal
embedded in vacuum is biased downwards at large $r$. Centres are drawn preferentially from
the domain interior; $r$ is expressed in units of the first-shell distance over
$[0.5, 4.0]$ in $100$ bins. A total $g(r)$ and the structure factor $\hat{S}(k)$ form the
remaining two channels.
\TODO{VERIFY: the exact $S(k)$ construction (explicit $k$-vectors versus Debye sum, and the
powder-averaging convention) in matcher\_v3\_validated\_structure\_factor/structure\_factor.py.}

\paragraph{Distance.} Curves are compared by the one-dimensional Wasserstein distance. Each
curve is clipped at zero and normalised to unit mass, so only the position of mass along the
abscissa is compared and amplitude is discarded. Because both distributions share a support,
the distance is evaluated in closed form,
\begin{equation}
  W_1(P,Q) \;=\; \int \bigl| F_P(x) - F_Q(x) \bigr| \, dx ,
  \label{eq:wasserstein}
\end{equation}
the area between the cumulative distributions. No multidimensional transport problem is
solved.

\paragraph{Ranking and acceptance.} Raw channel distances are not commensurate, so each is
divided by the median distance between two \emph{different} library anchors on that channel.
The ranking statistic is the mean of the scaled channel distances,
\begin{equation}
  D_f \;=\; \frac{1}{|C|}\sum_{c \in C} \frac{d_c(\mathcal{D}, f)}{\bar{d}_c},
  \label{eq:matcher-D}
\end{equation}
which is dimensionless and calibrated so that $D \approx 1$ corresponds to the typical
separation between two unrelated prototypes. Acceptance uses per-family, per-channel
tolerances $\tau_{f,c}$ estimated from the thermal ensemble as an inflated high quantile of
the envelope distances, combined conservatively,
\begin{equation}
  S_f \;=\; \max_{c} \frac{d_c(\mathcal{D}, f)}{\tau_{f,c}} ,
  \label{eq:matcher-S}
\end{equation}
so a good score on one channel cannot mask a failure on another. A candidate is matched to
$f^{\ast} = \arg\min_f D_f$ when $D_{f^{\ast}}$ is below the acceptance cutoff and the
runner-up margin exceeds a minimum; otherwise the candidate is reported as
library-unmatched or ambiguous. A stoichiometry gate on the LARGE/SMALL ratio is applied
before the structural channels.

The frozen constants are an acceptance cutoff of $0.40$, a minimum margin of $1.15$, an
envelope quantile of $0.90$ inflated by $1.30$, and a stoichiometry tolerance of $0.15$.
Note that $\hat{S}(k)$, the physical structure factor, and $S_f$, the matcher acceptance
score, are unrelated quantities that unfortunately share a symbol in the literature; they
are kept distinct here.

\subsection{Reference construction and a known limitation}
\label{sec:refbuild}

% SOURCE:
% crystal_discovery/cif_library/README.md
% crystal_discovery/crystal_stability/runs/*/configuration.yaml
Reference configurations are built from CIF prototypes by tiling the unit cell and choosing
a lattice scale that minimises the interaction energy of the colloidal potential. The
production builds optimised a \emph{single uniform} scale factor, which preserves the axial
ratios of the source CIF. For a cubic cell this is correct, since the cell has one free
length. For a non-cubic cell it is not: the axial ratios of the source describe an atomic
crystal, whereas the ideal ratios for a colloidal crystal of the same structure type depend
on the particle radius ratio.

Of the $42$ catalogued prototypes, $19$ are non-cubic. Re-optimising each
symmetry-distinct axis independently leaves the hexagonal Ni$_3$Sn reference unchanged but
alters the orthorhombic zangenite reference substantially ($c/a$ from $0.480$ to $0.596$).
Results that depend on the zangenite anchor are flagged accordingly in \cref{sec:results}.
All $122$ built references pass a cohesion check, with nearest-neighbour spacings between
$1.02$ and $1.11$ times the steric contact distance.

\subsection{Evaluation metrics}
\label{sec:metrics}

% SOURCE:
% results/motif_pipeline/core3_atlas_v3/scripts/v3_03_compare.py
Detection is evaluated by AUROC, average precision, true-positive rate at the calibrated
threshold and false-positive rate on known domains. Threshold stability is assessed by
resampling calibration \emph{groups} and reporting the resulting spread in $\tau$ together
with the fraction of validation domains whose accept/reject decision changes. Robustness to
the choice of descriptor centres is assessed by re-scoring each validation domain under
alternative centre sets. Matching is evaluated by top-1 reference, acceptance rate,
abstention rate under reference removal, and wrong-family acceptance.
